\subsection{Homotopisms of groupoids}

In this section, we shall use the notation \(u \sim v\) to express that two morphisms of groupoids u and v are homotopic.

DEFINITION 2. --- Let G, G' be groupoids and let \(\varphi\) be a morphism from G to G'. An inverse up to homotopy of \(\varphi\) is a morphism of groupoids \(\psi\) from G' to G such that the morphisms \(\psi \circ \varphi\) and \(\varphi \circ \psi\) are respectively homotopic to \(\mathrm{Id}_{\mathrm{G}}\) and to \(\mathrm{Id}_{\mathrm{G}'}\) . One says that \(\varphi\) is a homotopism if there exists an inverse of \(\varphi\) up to homotopy.

An isomorphism of groupoids is a homotopism.

Let G and G' be groupoids. Let \(\varphi\) , \(\varphi'\) be morphisms of groupoids from G to G' that are homotopic. If \(\varphi\) is a homotopism, then so is \(\varphi'\) . Indeed, if \(\psi\) denotes a homotopy inverse of \(\varphi\) , we have \(\psi \circ \varphi' \sim \psi \circ \varphi \sim \mathrm{Id}_{\mathrm{G}}\) and \(\varphi' \circ \psi \sim \varphi \circ \psi \sim \mathrm{Id}_{\mathrm{G}'}\) , which proves that \(\psi\) is a homotopy inverse of \(\varphi'\) .

Let G, G', and G'' be groupoids and let \(\varphi\colon\mathrm{G}\to\mathrm{G}'\) , \(\varphi'\colon\mathrm{G}'\to\mathrm{G}''\) , \(\psi\colon\mathrm{G}'\to\mathrm{G}\) , \(\psi'\colon\mathrm{G}''\to\mathrm{G}'\) be morphisms of groupoids. Then, among the following conditions:

\begin{enumerate}
\def\labelenumi{(\roman{enumi})}
\item
  \(\psi\) is a homotopy inverse of \(\varphi\) ;
\item
  \(\psi'\) is a homotopy inverse of \(\varphi'\) ;
\item
  \(\psi \circ \psi'\) is a homotopy inverse of \(\varphi' \circ \varphi\) ;
\end{enumerate}

any two imply the third. Indeed, suppose first that (i) and (ii) are satisfied; we then have

\[
\psi \circ \psi^ {\prime} \circ \varphi^ {\prime} \circ \varphi \sim \psi \circ \mathrm{Id} _ {\mathrm{G} ^ {\prime}} \circ \varphi \sim \psi \circ \varphi \sim \mathrm{Id} _ {\mathrm{G}}
\]

and, similarly, \(\varphi' \circ \varphi \circ \psi \circ \psi' \sim \mathrm{Id}_{\mathrm{G}'}\) , whence (iii). If (i) and (iii) are satisfied,

\[
\varphi^ {\prime} \circ \psi^ {\prime} \sim \varphi^ {\prime} \circ \varphi \circ \psi \circ \psi^ {\prime} \sim \mathrm{Id} _ {\mathrm{G} ^ {\prime \prime}}
\]

and

\[
\psi^ {\prime} \circ \varphi^ {\prime} \sim (\varphi \circ \psi) \circ \psi^ {\prime} \circ \varphi^ {\prime} \circ (\varphi \circ \psi) \sim \varphi \circ \psi \sim \mathrm{Id} _ {\mathrm{G} ^ {\prime}},
\]

whence condition (ii). The proof that conditions (ii) and (iii) imply condition (i) is analogous.

In particular, if two of the morphisms \(\varphi\) , \(\varphi'\) , \(\varphi'\circ\varphi\) are homotopisms, the same holds for the third.

PROPOSITION 1. --- Let G, G' be groupoids, let \(\varphi\) be a morphism from G to G', and let A be a subset of the vertex set of G that meets every orbit of G. For \(\varphi\) to be a homotopism, it is necessary and sufficient that the following conditions be satisfied:

\begin{enumerate}
\def\labelenumi{(\roman{enumi})}
\item
  the map \(\mathrm{Orb}(\varphi)\) from \(\mathrm{Orb}(\mathrm{G})\) to \(\mathrm{Orb}(\mathrm{G}')\) , deduced from \(\varphi\) by passing to orbits, is bijective;
\item
  for all \(a \in \mathrm{A}\) , the homomorphism \(\varphi_a \colon \mathrm{G}_a \to \mathrm{G}_{\varphi(a)}'\) is bijective.
\end{enumerate}

Suppose first that \(\varphi\) is a homotopism and let \(\psi\) be an inverse of \(\varphi\) up to homotopy. Then,

\[
\operatorname{Orb} (\psi) \circ \operatorname{Orb} (\varphi) = \operatorname{Orb} (\psi \circ \varphi) = \operatorname{Orb} (\operatorname{Id} _ {\mathrm{G}}) = \operatorname{Id} _ {\operatorname{Orb} (\mathrm{G})},
\]

because two homotopic groupoid morphisms induce the same map by passing to orbits. Likewise, \(\mathrm{Orb}(\varphi) \circ \mathrm{Orb}(\psi) = \mathrm{Id}_{\mathrm{Orb}(G')}\) . The map \(\mathrm{Orb}(\varphi)\) is therefore bijective, whence assertion (i). We also have, for every vertex a of G,

\[
\psi_ {\varphi (a)} \circ \varphi_ {a} = (\psi \circ \varphi) _ {a};
\]

since \(\psi \circ \varphi\) is homotopic to \(\mathrm{Id}_{\mathrm{G}}\) , the homomorphism \((\psi \circ \varphi)_a\) is bijective (cf. p. 180), so that \(\varphi_a\) is injective and \(\psi_{\varphi(a)}\) is surjective. Exchanging the roles of \(\varphi\) and \(\psi\) , we see that \(\varphi_a\) is also surjective, whence condition (ii).

Now suppose that conditions (i) and (ii) are satisfied and let us prove that \(\varphi\) is a homotopism.

First treat the case where each orbit of G' is reduced to a point.

For each vertex \(b\) of \(G'\) , choose a vertex \(u(b)\) of \(G\) belonging to \(A\) and whose image under \(\varphi\) is \(b\) . This is possible because the map \(\operatorname{Orb}(\varphi)\) is surjective and \(A\) meets each orbit of \(G\) . Let \(f\) be an arrow of \(G'\) ; we have \(o(f) = t(f)\) by hypothesis; set \(b = o(f)\) and \(a = u(b)\) . By condition (ii), there exists a unique arrow \(v(f) \in G_a\) whose image under \(\varphi\) is \(f\) . The pair \(\psi = (u, v)\) is a morphism of groupoids from \(G'\) to \(G\) . Let us prove that \(\psi\) is an inverse of \(\varphi\) up to homotopy. We already have \(\varphi \circ \psi = \mathrm{Id}_{\mathrm{G}'}\) by construction of \(\psi\) .

Let \(x\) be a vertex of G. Set \(a = \psi(\varphi(x))\) ; this is an element of A such that \(\varphi(a) = \varphi(x)\) . Since \(\mathrm{Orb}(\varphi)\) is injective, \(a\) belongs to the orbit of \(x\) in G and there exists an arrow \(f\) in G connecting \(a\) to \(x\) . The arrow \(h(x) = \psi(\varphi(f))^{-1}f\) then connects \(a\) to \(x\) and one has \(\varphi(h(x)) = e_{\varphi(a)}\) .

Let us show that the map \(h \colon \operatorname{Som}(\mathbf{G}) \to \operatorname{Fl}(\mathbf{G})\) thus defined is a homotopy connecting \(\psi \circ \varphi\) to \(\operatorname{Id}_{\mathbf{G}}\) . Condition (i) of definition 1 is satisfied, by construction. Let \(f\) be an arrow of \(\mathbf{G}\) , \(x\) its origin and \(y\) its target. We have \(\varphi(x) = \varphi(y)\) ; set \(a = \psi(\varphi(x)) = \psi(\varphi(y))\) . The arrows \(h(x)fh(y)^{-1}\) and \(\psi \circ \varphi(f)\) belong to \(\mathrm{G}_a\) and both have image \(\varphi(f)\) in \(\mathrm{G}'_{\varphi(a)}\) . As the map \(\varphi_a\) is injective, we have \(h(x)fh(y)^{-1} = \psi \circ \varphi(f)\) . Condition (ii) of definition 1 is thus satisfied.

Now let us prove proposition 1 in the general case. Let X be a maximal oriented forest of G' (II, p. 157, prop. 1). Let G'' be the groupoid deduced from G' by contracting the arrows of X and let \(\varphi'\colon G' \to G''\) be the canonical morphism. The morphism \(\varphi'\) satisfies the conditions of prop. 1 (II, p. 170, remark 1 and p. 178, cor. 2), and so does the morphism \(\varphi'\circ\varphi\) .

Since the orbits of G'' are reduced to points, it follows from the particular case already proved that \(\varphi'\) and \(\varphi'\circ\varphi\) are homotopisms, so that \(\varphi\) is also one. This completes the proof of the proposition.

COROLLARY 1. --- Let G be a groupoid, let A be a set, and let \(f\colon A \to \operatorname{Som}(G)\) be a map. If the image of \(f\) meets each orbit of G, the canonical groupoid morphism \(\varphi\) from \(f^{*}G\) to G is a homotopism.

By definition of the inverse image groupoid (II, p. 166, example 4), A is the set of vertices of the groupoid \(f^{*}\mathrm{G}\) and one has \(\mathrm{Fl}_{a,b}(f^{*}\mathrm{G}) = \mathrm{Fl}_{f(a),f(b)}(\mathrm{G})\) for every pair \((a,b)\) of elements of A. Moreover, we have \(\operatorname{Som}(\varphi) = f\) and \(\operatorname{Fl}(\varphi)\) induces the identity map from \(\operatorname{Fl}_{a,b}(f^{*}\mathrm{G})\) to \(\operatorname{Fl}_{f(a),f(b)}(\mathrm{G})\) . Consequently, the map \(\operatorname{Orb}(\varphi)\) is bijective and the homomorphism \(\varphi_{a}\colon (f^{*}\mathrm{G})_{a} \to \operatorname{G}_{f(a)}\) is an isomorphism, for all \(a \in A\) . The hypotheses of Proposition 1 are therefore satisfied.

COROLLARY 2. --- Let G be a groupoid, let X be an oriented forest of G, and let G' be the groupoid obtained from G by contracting the arrows of X. The canonical morphism from G to G' is a homotopism.

This follows from Remark 1 of II, p. 170, and Corollary 2 of II, p. 178, which show that the hypotheses of Proposition 1 are satisfied.

